\clearpage
\unitheader{E}{Optional exploration}{}

To explore the topics in this assignment more broadly, construct your own prediction
problem. This optional exercise can lead to bonus points on the quiz, as described
in the Instructions.

\begin{problem}{Your own prediction problem (optional)}
You may have found an unexpected interaction between schedules, LR and WD, batch
size, or model size. Use one such observation to design a question that a classmate
could reason about from the assignment, but whose answer goes beyond the supplied examples.

Include the following in your exploration:
\begin{parts}
\item A question writeup in the style of the examples. Provide answer options and
state whether to select one or all that apply. Record your prediction and reasoning
before testing the target intervention. Make the outcome objectively checkable from
an experiment with specified model, data, training budget, and evaluation metric.

\item A compact code diff (under 100 lines), executable code, and the measured evidence
that determines the answer. Include the answer key and explain any difference from
your prediction. Keep the diff focused on the intervention and reuse the shared training
interface. If variation across runs could change the correct option, check it and revise
the question or option ranges accordingly.
\end{parts}

Aim for a fair, novel question with plausible alternatives. Start from an experiment
you already understand and use the smallest informative comparison. Course staff may
use interesting questions in a future quiz. This exploration is additional to the
starting grids in the four parts.
\end{problem}
